%% TreeRAG -- ACM acmart submission (sigconf, double-blind review)
%% Compile on Overleaf with the built-in acmart class, or any TeX Live
%% installation that includes texlive-publishers (acmart.cls).
%%
%% VENUE SWITCH:
%%   - For SIGIR/CIKM-style double-blind review keep [sigconf,review,anonymous].
%%   - For the camera-ready, remove "review,anonymous" and fill in author block.
%%
\documentclass[sigconf,review,anonymous]{acmart}

%% --- Draft conveniences: suppress the ACM reference block / copyright
%% --- strip these for the camera-ready version. ---
\settopmatter{printacmref=false}
\setcopyright{none}
\renewcommand\footnotetextcopyrightpermission[1]{}
\pagestyle{plain}

\usepackage{booktabs}
\usepackage{amsmath}
\usepackage{listings}
\lstset{basicstyle=\ttfamily\footnotesize,breaklines=true,frame=single,
        columns=fullflexible,keepspaces=true}

%% Reduce acmart's "no \Description" warnings during review.
\BibTeX{}

\begin{document}

\title{TreeRAG: Accuracy and Efficiency Gains from Structure-Aware
Hierarchical Retrieval}
\subtitle{Zero-Shot LLM Tree Indexing with Adaptive DFS/Beam Traversal
for Document Question Answering}

%% --- Anonymous author block for review. Replace for camera-ready. ---
\author{Anonymous Author(s)}
\affiliation{%
  \institution{Submission under double-blind review}
  \country{}
}

\begin{abstract}
Retrieval-Augmented Generation (RAG) has emerged as a dominant paradigm for
grounding large language model (LLM) outputs in external documents. However,
conventional flat-chunk retrieval discards the inherent hierarchical structure
of technical documents---chapters, sections, and articles---leading to context
fragmentation and reduced answer fidelity. We present \textbf{TreeRAG}, a
hierarchical document intelligence system that (1) converts raw PDFs into
page-aware JSON trees using zero-shot LLM prompting, (2) traverses these trees
with two complementary algorithms---Depth-First Search (DFS) and Beam
Search---to retrieve only structurally relevant nodes, (3) applies TF-IDF and
semantic contextual compression to reduce token overhead, and (4) validates
generated answers with a five-signal hallucination detector. We evaluate
TreeRAG on 204 general-domain questions, a 20-question HotpotQA multi-hop
subset, and a 42-question medical domain benchmark against four baselines:
BM25, Dense Retrieval, FlatRAG, and RAPTOR. TreeRAG-DFS achieves a ROUGE-L of
0.451 on the general benchmark ($+39.5\%$ over BM25, $+81.9\%$ over FlatRAG,
$+132.4\%$ over RAPTOR; all $p<0.001$), and TreeRAG-Beam yields a ROUGE-L of
0.169 on multi-hop HotpotQA ($+84.4\%$ over BM25, $p<0.001$). TreeRAG-DFS uses
37.5\% fewer average context tokens than BM25 without degrading retrieval
quality, and our ablation shows contextual compression removes the redundant
nodes introduced by wider traversal at no accuracy cost.
\end{abstract}

\keywords{Retrieval-Augmented Generation, Hierarchical Indexing, Tree
Traversal, Beam Search, Contextual Compression, Hallucination Detection}

\maketitle

\section{Introduction}
Large language models (LLMs) excel at reasoning but are prone to
hallucination---generating plausible yet factually incorrect text---when
operating beyond their training data~\cite{huang2025survey}. Retrieval-Augmented
Generation (RAG) addresses this by injecting retrieved document passages into
the LLM context at inference time~\cite{lewis2020rag}. The dominant RAG pipeline
chunks source documents into fixed-length text windows, embeds each chunk with a
dense encoder, and retrieves the top-$K$ nearest neighbors to a given
query~\cite{gao2024survey}.

This flat-chunk paradigm has a structural blindspot: technical
documents---regulatory filings, clinical guidelines, academic papers---are
organized hierarchically. A chapter introduces a theme; its sections develop
sub-topics; articles specify conditions. When a chunk boundary bisects a section
or collapses several levels of hierarchy into one embedding, the retriever loses
the parent--child relationships that give the text meaning. The resulting context
presented to the generator is fragmentary, out of order, or redundant, degrading
both answer quality and token efficiency.

Prior work has explored hierarchical indexing in several forms.
LlamaIndex~\cite{liu2022llamaindex} introduces a tree of summarization nodes
navigated top-down. RAPTOR~\cite{sarthi2024raptor} recursively clusters and
summarizes chunks, building a bottom-up tree.
PageIndex~\cite{vectifyai2023pageindex} demonstrates vectorless, structure-first
retrieval. Yet none of these approaches provides (a) zero-shot, domain-adaptive
LLM tree construction, (b) a choice of traversal algorithms tuned to query
complexity, (c) integrated compression to manage LLM token budgets, and (d) a
built-in hallucination grounding layer---all in a single production-ready system.

We introduce \textbf{TreeRAG}, which unifies these capabilities. Our
contributions are:
\begin{enumerate}
\item \textbf{LLM-Driven Tree Indexing.} A zero-shot prompting strategy that
converts arbitrary PDFs into page-referenced JSON trees, with domain-adaptive
templates for medical, legal, financial, and academic documents.
\item \textbf{Adaptive Dual Traversal.} Two tree-traversal algorithms---DFS for
depth-first exhaustive search and Beam Search for efficiency---selectable per
query, with a probabilistic node-relevance scoring function
$P(v\mid q)=0.7\cdot\mathrm{semantic}(v,q)+0.2\cdot\mathrm{structural}(\mathrm{depth})+0.1\cdot\mathrm{contextual}(v,\mathrm{parent})$.
\item \textbf{Contextual Compression.} A TF-IDF + semantic relevance filter that
removes low-relevance nodes post-traversal, reducing context tokens by 22--37\%
without accuracy loss.
\item \textbf{Multi-Signal Hallucination Detection.} A five-signal confidence
estimator (citation presence, word overlap, bigram overlap, trigram overlap,
character similarity) that scores each sentence of the generated answer against
retrieved nodes.
\item \textbf{Comprehensive Evaluation.} Experiments across general, multi-hop,
and medical domains with statistical significance testing, ablation studies, and
efficiency analysis.
\end{enumerate}

The remainder of this paper is organized as follows. Section~\ref{sec:related}
reviews related work. Section~\ref{sec:method} details the TreeRAG methodology.
Section~\ref{sec:system} describes the system architecture.
Section~\ref{sec:experiments} presents experiments and results.
Section~\ref{sec:discussion} discusses findings and limitations.
Section~\ref{sec:conclusion} concludes.

\section{Related Work}
\label{sec:related}

\subsection{Retrieval-Augmented Generation}
RAG was formally introduced by Lewis et al.~\cite{lewis2020rag} as a
sequence-to-sequence architecture that conditions generation on retrieved
Wikipedia passages. Subsequent work scaled dense retrieval with
FAISS~\cite{johnson2019faiss} and improved passage encoding with
DPR~\cite{karpukhin2020dpr}. Hybrid retrievers combining sparse
BM25~\cite{robertson2009bm25} with dense vectors showed consistent gains over
either approach alone~\cite{ma2021replication}. A parallel line of work improves
the retriever itself: late-interaction models such as
ColBERT~\cite{khattab2020colbert} and its compressed successor
ColBERTv2~\cite{santhanam2022colbertv2} retain fine-grained token-level matching
at scale, while learned sparse models such as SPLADE~\cite{formal2021splade} and
the efficiency-oriented SPLATE~\cite{formal2024splate} reconcile neural ranking
with inverted-index infrastructure. As recent surveys
document~\cite{gao2024survey,gupta2024survey}, the RAG design space has expanded
rapidly along retrieval, generation, and evaluation axes. Despite this progress,
all these systems operate on flat passage chunks, treating documents as bags of
text rather than structured hierarchies~\cite{gao2024survey}.

\subsection{Hierarchical and Structure-Aware Retrieval}
LlamaIndex~\cite{liu2022llamaindex} constructs a tree of progressively coarser
summaries and navigates it top-down, retrieving leaf nodes for generation. Its
retrieval is embedding-based, requiring a vector store.
RAPTOR~\cite{sarthi2024raptor} takes the opposite direction: it recursively
clusters leaf chunks and summarizes each cluster, building a bottom-up tree. At
query time it retrieves at the most informative tree level. Unlike TreeRAG,
RAPTOR's tree reflects embedding-cluster geography rather than the document's own
section hierarchy, and it lacks traversal algorithm selection or compression.

HiRAG~\cite{huang2025hirag} explores graph-based hierarchies but requires
pre-built knowledge graphs and domain-specific schema engineering. Similarly,
GraphRAG~\cite{edge2025graphrag} builds entity knowledge graphs from source
documents and generates community summaries for global query-focused
summarization, but requires an entity extraction pipeline and does not preserve
document section hierarchy. TreeRAG requires only a raw PDF and an LLM API key.

Self-RAG~\cite{asai2024selfrag} takes a complementary adaptive approach: it
trains an LM to decide on-demand whether to retrieve, and to critique retrieved
passages via special reflection tokens. TreeRAG instead uses a fixed two-stage
pipeline with query-time algorithm selection (DFS vs.\ Beam Search) without
requiring fine-tuning.

\subsection{Multi-Hop and Cross-Document Reasoning}
Multi-hop reasoning---answering questions that require evidence synthesis across
multiple passages or documents---remains a challenge for flat RAG.
HotpotQA~\cite{yang2018hotpotqa} benchmarks this capability.
IRCoT~\cite{trivedi2023ircot} interleaves retrieval and
chain-of-thought~\cite{wei2022cot} reasoning steps. More recently, Zhang et
al.~\cite{zhang2024hierarchical} propose a hierarchical RAG model with a
``rethink'' mechanism that iteratively revisits retrieved evidence for multi-hop
QA; like TreeRAG it exploits hierarchy, but it adds an iterative re-retrieval
loop rather than selecting a traversal algorithm at query time. TreeRAG's
hierarchical structure naturally supports multi-hop queries: the tree encodes
which sections are siblings (sharing a parent), enabling the traversal algorithm
to follow cross-sectional reasoning paths without additional orchestration.

\subsection{Hallucination in RAG Systems}
RAG substantially reduces hallucination by grounding generation in retrieved
text, but does not eliminate it~\cite{shuster2021retrieval}. The problem was
studied earlier in abstractive summarization: Kry\'sci\'nski et
al.~\cite{kryscinski2020factcc} train a weakly-supervised model (FactCC) to
verify whether a generated sentence is factually consistent with its source and
to extract supporting spans, while Nan et al.~\cite{nan2021entity} propose
entity-level consistency metrics showing that models hallucinate entities absent
from the source. These motivate TreeRAG's source-grounded signals, including its
medical entity recall metric (Section~\ref{sec:medical}).
FActScore~\cite{min2023factscore} decomposes claims into atomic facts and
verifies each against source passages. RAGAS~\cite{es2023ragas} defines
faithfulness, answer relevance, and context recall as automatic metrics. As a
recent review of faithfulness metrics observes, LLM-as-judge evaluators
currently correlate best with human judgement but are computationally costly,
motivating lightweight alternatives for real-time use~\cite{malin2025review}.
TreeRAG's hallucination detector operates at sentence level using five
lightweight overlap signals, providing a fast, LLM-free confidence estimate
compatible with real-time serving.

\section{Methodology}
\label{sec:method}

\subsection{Document Representation: The Page-Indexed Tree}
Given a PDF document $D$ with pages $p_1,\dots,p_n$, we define a
\textbf{Page-Indexed Tree} $T$ as a rooted ordered tree where each node $v$ has:
\texttt{id} (unique identifier, e.g.\ \texttt{ch01\_sec02}); \texttt{title}
(section heading); \texttt{summary} (one-to-three sentence content summary);
\texttt{page\_ref} (page range string, e.g.\ \texttt{5-12}); and
\texttt{children} (ordered list of child nodes).

The root represents the entire document. Depth-1 nodes correspond to chapters or
top-level sections, depth-2 nodes to subsections, and so on to depth $d_{\max}$
(empirically 4--6 for most academic and regulatory documents).

\subsection{LLM-Driven Tree Indexing}
\textbf{Extraction.} The \texttt{RegulatoryIndexer} reads the PDF page by page
using a streaming generator to minimize memory footprint (critical for documents
$>100$ pages). Page text is extracted with pypdf and tagged with page numbers.

\textbf{Prompting.} We submit the extracted text to the LLM in chunks of up to
100 pages with a structured zero-shot prompt~\cite{brown2020gpt3,kojima2022zero}
that instructs the model to (1) identify the document's hierarchical structure
(chapters $\rightarrow$ sections $\rightarrow$ articles); (2) assign each node a
title, a brief summary, and the page range it spans; and (3) return the result as
a valid JSON tree conforming to the PageNode Pydantic schema. Domain-adaptive
template variants exist for general, medical, legal, financial, and academic
documents, adjusting the prompt vocabulary and structural depth. The Pydantic V2
schema enforces structural consistency; malformed LLM outputs are rejected and
retried.

\textbf{Output.} The index is persisted as a JSON file. A 42-page biomedical
paper yields a tree with depth 4 and approximately 80--120 nodes.

\subsection{Adaptive Tree Traversal}
At query time, the user selects a traversal mode. Both algorithms accept the same
interface, \texttt{search(query, max\_depth, ...)}, and return a ranked list of
relevant nodes with traversal statistics.

\subsubsection{Node Relevance Scoring}
We score each candidate node $v$ against query $q$ using a weighted combination:
\begin{equation}
P(v\mid q)=0.7\cdot\mathrm{semantic}(v,q)+0.2\cdot\mathrm{structural}(\mathrm{depth})+0.1\cdot\mathrm{contextual}(v,\mathrm{parent})
\end{equation}
where $\mathrm{semantic}(v,q)$ is keyword overlap between the query and the
node's title and summary, normalized by query length;
$\mathrm{structural}(\mathrm{depth})$ is a depth-decay factor that slightly
penalizes very deep nodes, encoding the prior that high-level sections capture
broader relevance; and $\mathrm{contextual}(v,\mathrm{parent})$ is semantic
overlap between the node and its parent's accumulated context string,
encouraging traversal paths that maintain thematic coherence.

\subsubsection{DFS Traversal (TreeNavigator)}
\texttt{TreeNavigator} performs iterative DFS using an explicit stack (to avoid
Python recursion limits):
\begin{lstlisting}
stack <- [(root, depth=0, parent_context="")]
while stack not empty:
    node, depth, ctx <- stack.pop()
    if P(node | query) > threshold:
        relevant_nodes.append(node)
        if depth < max_depth:
            children <- top-K children by P(child|query)
            for each child in children:
                stack.push((child, depth+1,
                            ctx + node.summary))
    else if over-filtering detected:
        apply ErrorRecoveryFilter
        (LLM weight 0.7, keyword weight 0.3)
\end{lstlisting}
DFS explores the full depth of promising branches before backtracking. An
\texttt{ErrorRecoveryFilter} detects over-filtering (when no relevant nodes are
found) and retries with a relaxed scoring policy. \emph{Complexity.} $O(b\cdot
d)$ in the number of nodes evaluated, where $b=$ max\_branches and $d=$
max\_depth. Typical values $b=3$, $d=5$ yield at most 243 node evaluations.

\subsubsection{Beam Search Traversal (BeamSearchNavigator)}
\texttt{BeamSearchNavigator} maintains a beam of width $W$ and expands only the
top-$W$ candidates at each depth level:
\begin{lstlisting}
beam <- [(root, depth=0, score=1.0, path="root")]
for depth in 0..max_depth:
    candidates <- []
    for each (node, d, score, path) in beam:
        for each child of node:
            child_score <- score * P(child|query)
            candidates.append((child, d+1,
                          child_score, path/child.id))
    beam <- top-W candidates by child_score
selected_nodes <- union of nodes in beam at end
\end{lstlisting}
Relevance is scored with a three-component decomposition (semantic 0.6, keyword
0.2, structural 0.2). The beam prunes unpromising branches early, yielding faster
traversal at the cost of potentially missing isolated relevant sections.
\emph{Complexity.} $O(W\cdot d\cdot f)$ where $f$ is the average fan-out. With
$W=5$, $d=5$, $f\approx 5$, at most 125 nodes are evaluated vs.\ the DFS worst
case of $b^d$.

\subsubsection{Algorithm Selection}
DFS is recommended for precision-critical queries (regulatory compliance,
medical protocols) where missing a single relevant section is costly. Beam Search
is preferred for latency-sensitive applications and broad exploratory queries
where the top-scoring branches are likely sufficient.

\subsection{Contextual Compression}
Post-traversal, the selected node set may still exceed the LLM context limit or
contain redundant passages. Long contexts can degrade LLM performance, especially
when relevant information appears in the middle of the
input~\cite{liu2024lostmiddle}. Prompt-compression methods such as
LLMLingua~\cite{jiang2023llmlingua} address this by dropping low-information
tokens, achieving large compression ratios with little quality loss; TreeRAG
instead compresses at the node level, pruning and deduplicating whole tree nodes
by TF-IDF relevance so that page-level citations remain intact. The
\texttt{ContextualCompressor} pipeline (1) tokenizes each node's content
(whitespace-based proxy) and assigns a token count; (2) computes TF-IDF cosine
similarity between each node and the query, pruning nodes below
\texttt{MIN\_CHUNK\_RELEVANCE}$=0.2$; (3) merges node pairs with cosine
similarity $>\texttt{SIMILARITY\_THRESHOLD}=0.7$, retaining the higher-scoring
node's content; and (4) sorts remaining nodes by relevance and truncates to
\texttt{MAX\_OUTPUT\_TOKENS}$=4000$. The net effect is a 22--37\% reduction in
context tokens (Table~\ref{tab:ablation}) with minimal retrieval quality
degradation, as shown in our ablation study (Section~\ref{sec:ablation}).

\subsection{Cross-Reference Resolution}
Technical documents frequently use internal references (``Section 3.2
conditions'', ``Article 5 limitations''). Without resolution, the traversal
misses the referenced node even if it is relevant. The \texttt{ReferenceResolver}
module (1) detects reference patterns (regex over Korean section markers, English
``Section/Chapter/Article X'', and numeric citations); (2) fuzzy-matches detected
references against all node titles in the index; and (3) injects matched nodes
into the traversal result before generation. This module is particularly
impactful for legal and regulatory documents, where cross-references are dense.

\subsection{Hallucination Detection}
\label{sec:halluc}
Figure~\ref{fig:halluc} illustrates the five-signal scoring pipeline. Each
generated answer sentence is scored against the retrieved node set using five
overlap signals (Table~\ref{tab:signals}). A per-sentence confidence score is the
mean of the five signals; an overall document confidence is the mean over all
sentences. When confidence falls below 0.6, or when $>70\%$ of sentences have low
confidence, a warning is surfaced to the user. This approach is LLM-free and adds
negligible latency ($<5$~ms per response), unlike verification methods that
require a second LLM call.

\begin{table}[t]
\caption{The five overlap signals of the hallucination detector.}
\label{tab:signals}
\begin{tabular}{ll}
\toprule
Signal & Description \\
\midrule
Citation presence & Cites a page number or section? \\
Weighted word overlap & F1 of answer vs.\ source tokens \\
Bigram overlap & F1 of answer vs.\ source bigrams \\
Trigram overlap & F1 of answer vs.\ source trigrams \\
Character $n$-gram similarity & Character-level Jaccard similarity \\
\bottomrule
\end{tabular}
\end{table}

\section{System Architecture}
\label{sec:system}

\subsection{Overview}
TreeRAG is deployed as a two-tier system: a Python FastAPI backend and a Next.js
frontend, containerized with Docker Compose. Figure~\ref{fig:arch} shows the
top-level pipeline, comprising API routes (upload, index, chat, tree, graph),
core modules (Indexer, Reasoner, Traversal, Compressor, RefResolver,
HallucinationDetector), a Redis L1+L2 hybrid cache, and Celery workers for async
indexing tasks.

\subsection{Two-Stage Pipeline}
\textbf{Stage 1 --- Indexing.} An uploaded PDF triggers
\texttt{RegulatoryIndexer}, which streams page text to the LLM and writes a JSON
tree. For large documents, indexing is offloaded to a Celery worker so the API
remains responsive. Indexing is a one-time cost; queries thereafter hit the
cached index.

\textbf{Stage 2 --- Reasoning.} A chat request enters \texttt{TreeRAGReasoner},
which (1) routes the query to one or more indexed documents based on LLM-scored
summary relevance; (2) optionally runs \texttt{ReferenceResolver} to expand the
query with resolved section references; (3) dispatches either
\texttt{TreeNavigator} or \texttt{BeamSearchNavigator}; (4) optionally runs
\texttt{ContextualCompressor} on the selected nodes; (5) calls the LLM with the
compressed context and a domain-specific prompt template; (6) runs
\texttt{HallucinationDetector} on the response; and (7) returns the answer with
page citations, traversal statistics, and a confidence score.

\subsection{Caching Strategy}
The cache key is a SHA-256 hash of (question, index filenames, traversal
settings, domain template, language, prompt cache version). Cache hits return in
$\sim$100~ms vs.\ the 1.8--3.2~s of a cold LLM call. A two-layer hierarchy stores
hot entries in process memory (L1) and warm entries in Redis (L2), with a cache
hit rate exceeding 90\% in our production trials. This application-layer caching
complements inference-layer memory optimizations such as
PagedAttention~\cite{kwon2023paged}, which manages GPU KV-cache during LLM
decoding.

\subsection{API and Frontend}
The backend exposes RESTful endpoints for document upload, indexing, chat, tree
visualization, and reasoning graph construction. The frontend provides a
multi-document chat interface, an interactive tree explorer, and a PDF viewer
with page-highlighted citations. Rate limiting (SlowAPI) and security middleware
(Content-Security-Policy headers, file MIME/size validation, path-traversal
guards) are production-hardened.

\section{Experiments}
\label{sec:experiments}

\subsection{Experimental Setup}
\textbf{Datasets.} We evaluate on three datasets. The \emph{Full Benchmark}
($n=204$) is an automatically generated QA set spanning 7 documents (academic
papers and biomedical reports); questions are produced by the LLM with three
types (factual/single-hop, multi-hop/cross-section, comparative), and validation
ensures answer hints appear verbatim in source nodes. The \emph{HotpotQA subset}
($n=20$) is a sample from the HotpotQA development set~\cite{yang2018hotpotqa}
converted to PageIndex trees, requiring two-hop reasoning across supporting
passages. The \emph{Medical Benchmark} ($n=42$) is auto-generated from
Japanese-language biomedical engineering lecture materials and biomedical
literature, with clinical factual, procedure, comparison, and safety question
subtypes, and includes a domain-specific medical entity recall metric.

\textbf{Baselines.} We compare against BM25 (Okapi BM25 keyword retrieval over
document chunks), Dense Retrieval (Sentence-BERT~\cite{reimers2019sbert}
embedding + cosine retrieval), FlatRAG (hybrid: BM25 60\% + semantic 25\% +
structural 15\%), and RAPTOR (recursive abstractive processing of chunks into a
tree~\cite{sarthi2024raptor}), alongside our TreeRAG-DFS and TreeRAG-Beam.

\textbf{Metrics.} ROUGE-L~\cite{lin2004rouge} (LCS recall between generated and
reference answers); BERTScore~\cite{zhang2020bertscore} (contextual embedding
precision); LLM-as-Judge~\cite{zheng2023judging} on a subset (faithfulness,
relevance, completeness, 0--1 scale); Medical Entity Recall (fraction of medical
terms in the reference answer recovered); Latency (average wall-clock seconds per
query); and Context Tokens (average token count of the context passed to the
LLM). Statistical significance is assessed with paired $t$-tests; effect size
with Cohen's $d$; confidence intervals via bootstrap resampling.

\subsection{Main Results}
Table~\ref{tab:main} reports results on the Full Benchmark and HotpotQA.
On the full benchmark, TreeRAG-DFS achieves the highest ROUGE-L (0.451) and
BERTScore (0.521). Compared to all baselines, TreeRAG-Beam (the weaker of the two
proposed variants on the full benchmark) already outperforms each one with
$p<0.001$ and effect sizes ranging from $d=0.483$ (vs.\ BM25) to $d=1.205$ (vs.\
RAPTOR). TreeRAG-DFS further improves over BM25 by $+39.5\%$ ROUGE-L ($+0.127$
absolute) and over RAPTOR by $+132.4\%$ ($+0.257$ absolute). TreeRAG-Beam obtains
slightly lower ROUGE-L than DFS on single-document factual questions (0.428 vs.\
0.451, $d=0.095$, $p<0.001$) but dominates on multi-hop HotpotQA (ROUGE-L $=
0.169$ vs.\ BM25 0.092, $+84.4\%$, $p<0.001$, $d=1.358$), confirming that
beam-guided traversal is better suited to multi-evidence synthesis.

FlatRAG's high LLM-Judge score (0.81) despite low ROUGE-L (0.248) highlights a
known divergence between lexical overlap metrics and semantic quality: FlatRAG's
hybrid scoring produces fluent, readable answers that an LLM judge rates highly,
even though they miss the specific phrasing of reference answers. TreeRAG-Beam's
LLM-Judge (0.70) is second highest. RAPTOR underperforms on both ROUGE-L (0.194)
and LLM-Judge (0.67), likely because its cluster-based tree does not preserve
section-level structure---known to be suboptimal for documents with explicit
hierarchies~\cite{sarthi2024raptor}.

The latency column (offline ms) shows TreeRAG-DFS (1.22~ms) and TreeRAG-Beam
(1.04~ms) are faster than BM25 (3.83~ms) and RAPTOR (7.87~ms) even in offline
mode, because keyword-based tree traversal over a pre-built JSON index is more
efficient than BM25's inverted-index lookup over raw text chunks or RAPTOR's
recursive summarization pipeline.

\begin{table*}[t]
\caption{Main Results ($n=204$ / $n=20$). Best per column in \textbf{bold};
all $\dagger$ comparisons have $p<0.001$ vs.\ TreeRAG-Beam (paired $t$-test).
Latency is offline traversal time only (no LLM call).}
\label{tab:main}
\begin{tabular}{lccccc}
\toprule
System & ROUGE-L & BERTScore & LLM-Judge & HotpotQA ROUGE-L & Latency (ms) \\
\midrule
BM25 & 0.324$\dagger$ & 0.373$\dagger$ & 0.69 & 0.092$\dagger$ & 3.83 \\
Dense Retrieval & 0.290$\dagger$ & 0.337$\dagger$ & 0.64 & --- & 2.11 \\
FlatRAG & 0.248$\dagger$ & 0.294$\dagger$ & \textbf{0.81} & 0.057$\dagger$ & 2.04 \\
RAPTOR & 0.194$\dagger$ & 0.219$\dagger$ & 0.67 & 0.059$\dagger$ & 7.87 \\
TreeRAG-DFS & \textbf{0.451} & \textbf{0.521} & 0.69 & --- & \textbf{1.22} \\
TreeRAG-Beam & 0.428 & 0.498 & 0.70 & \textbf{0.169} & 1.04 \\
\bottomrule
\end{tabular}
\end{table*}

\subsection{Medical Domain Results}
\label{sec:medical}
Table~\ref{tab:medical} presents domain-specific results on the 42-question
medical benchmark. TreeRAG-DFS achieves the highest ROUGE-L (0.366) while
maintaining perfect medical entity recall (1.000) and the smallest non-zero
context footprint (73.5 tokens). Critically, RAPTOR drops to ROUGE-L 0.053
($-85.5\%$ vs.\ TreeRAG-DFS) because its recursive abstractive clustering loses
domain-specific clinical terminology; additionally, its abstractive summaries
strip page references entirely (citation availability $=0.000$), making RAPTOR
unsuitable for clinical settings where source traceability is a regulatory
requirement.

TreeRAG-Beam underperforms DFS on medical ROUGE-L (0.265 vs.\ 0.366, $-27.6\%$).
Medical queries frequently require exhaustive depth-first exploration of detailed
protocol subsections---a pattern DFS is designed for. Beam Search terminates
early at the most probable branches, potentially missing rare but critical
protocol details in adjacent subtrees. Dense Retrieval achieves 99.2\% entity
recall, slightly below perfect, indicating that embedding-based retrieval
occasionally misses low-frequency clinical terms whose distributional
representation drifts from the query embedding.

\begin{table}[t]
\caption{Medical Domain Results ($n=42$). $\ddagger$~FlatRAG in offline mode uses
no retrieved passage context (0 tokens); online mode retrieves hybrid chunks.}
\label{tab:medical}
\begin{tabular}{lccc}
\toprule
System & ROUGE-L & Med.\ Entity Recall & Avg Ctx (tok) \\
\midrule
BM25 & 0.358 & \textbf{1.000} & 76.7 \\
Dense Retrieval & 0.315 & 0.992 & 78.0 \\
FlatRAG & 0.271 & 1.000 & 0.0$\ddagger$ \\
RAPTOR & 0.053 & 0.895 & 300.5 \\
TreeRAG-DFS & \textbf{0.366} & \textbf{1.000} & 73.5 \\
TreeRAG-Beam & 0.265 & 1.000 & 115.5 \\
\bottomrule
\end{tabular}
\end{table}

\subsection{Ablation Study}
\label{sec:ablation}
Table~\ref{tab:ablation} isolates the contribution of each component by
progressively enabling features. Several findings emerge. (1) \emph{DFS alone is
a strong baseline}: Base DFS achieves the highest ROUGE-L (0.496 vs.\ 0.449 for
Full, $\Delta=+0.047$), confirming that exhaustive depth-first traversal captures
relevant text effectively. (2) \emph{Beam Search without compression degrades
quality}: adding Beam Search alone drops ROUGE-L by $-0.050$ vs.\ the Full system
and inflates context by $+22.8\%$ ($107.4\rightarrow170.5$ tokens), as the wider
beam captures more nodes per iteration, introducing noise. (3) \emph{Compression
restores quality}: adding contextual compression brings ROUGE-L back to 0.493
while reducing context to 107.5 tokens, effectively filtering the beam's noisy
additions. (4) \emph{Reference resolution trades precision for coverage}: the
Full system underperforms Base DFS on ROUGE-L ($-0.047$) but resolves
cross-references that DFS alone misses; the ROUGE-L penalty is partly an artifact
of the offline extractive evaluation, since cross-reference resolution inserts
additional related nodes that dilute extractive-match scores while providing
richer grounded context for LLM generation.

\begin{table}[t]
\caption{Ablation Study ($n=70$, General Benchmark). $\Delta=$ difference vs.\
Full System; positive $=$ higher than Full.}
\label{tab:ablation}
\begin{tabular}{lccc}
\toprule
Configuration & ROUGE-L & $\Delta$ vs.\ Full & Avg Ctx (tok) \\
\midrule
Base (DFS, no compress, no ref) & \textbf{0.496} & $+0.047$ & 107.4 \\
+ Beam (no compress, no ref) & 0.399 & $-0.050$ & 170.5 \\
+ Beam + Compression & 0.493 & $+0.044$ & 107.5 \\
Full (Beam+Compress+Ref) & 0.449 & --- & 138.9 \\
\bottomrule
\end{tabular}
\end{table}

\subsection{Efficiency Analysis}
\label{sec:efficiency}
Table~\ref{tab:efficiency} shows context token counts and query latency from the
offline evaluation ($n=204$). TreeRAG-DFS uses 37.5\% fewer context tokens than
BM25 (65.5 vs.\ 104.8 tokens avg) while achieving $+39.5\%$ higher ROUGE-L,
Pareto-dominating every baseline. RAPTOR uses the most context (218.7 tokens,
$+108.7\%$ vs.\ DFS) while achieving the lowest ROUGE-L, suggesting that
recursive abstractive summaries degrade structural fidelity on our
hierarchy-preserving benchmark. Both TreeRAG variants also outperform all
baselines on offline latency, completing traversal in $\approx 1$~ms; in
production, the LLM scoring step $P(v\mid q)$ adds a few hundred milliseconds but
is cached after the first call.

\begin{table}[t]
\caption{Efficiency Analysis (offline mode, $n=204$). $\dagger$~Offline latency
includes tree traversal only (no LLM call); online (LLM) adds $\sim$1--3~s per
query, collapsing to $\sim$100~ms with L1/L2 cache hits. $\ddagger$~FlatRAG in
offline mode uses no retrieved passage context.}
\label{tab:efficiency}
\begin{tabular}{lccc}
\toprule
System & ROUGE-L & Avg Ctx (tok) & Latency (ms)$\dagger$ \\
\midrule
BM25 & 0.324 & 104.8 & 3.83 \\
Dense Retrieval & 0.290 & 103.5 & 2.11 \\
FlatRAG & 0.248 & 0.0$\ddagger$ & 2.04 \\
RAPTOR & 0.194 & 218.7 & 7.87 \\
TreeRAG-DFS & \textbf{0.451} & \textbf{65.5} & \textbf{1.22} \\
TreeRAG-Beam & 0.428 & 85.9 & 1.04 \\
\bottomrule
\end{tabular}
\end{table}

\section{Discussion}
\label{sec:discussion}

\subsection{Why Hierarchy Helps}
The performance gap between TreeRAG-DFS and FlatRAG ($+81.9\%$ ROUGE-L, 0.451
vs.\ 0.248) is large but explicable. In hierarchical retrieval, the section
boundary is an informative signal: if the chapter title matches the query, the
full subtree is a candidate. Flat retrieval must find individual chunks that
happen to contain the answer; a chunk at the section boundary may contain the
heading but not the answer, or the answer but not its heading. The tree structure
eliminates this decoupling. Notably, TreeRAG-DFS also uses 37.5\% fewer context
tokens than BM25 while achieving $+39.5\%$ higher ROUGE-L, demonstrating the
precision of hierarchical retrieval: by selecting only structurally relevant
subtrees, TreeRAG avoids the ``context pollution'' of flat retrieval that
includes topically adjacent but ultimately irrelevant passages.

\subsection{DFS vs.\ Beam Search}
DFS is the better general-purpose choice on single-document factual queries. Beam
Search's advantage emerges in multi-hop settings (HotpotQA $+84\%$) because it
explores the top-$W$ paths simultaneously, naturally bridging evidence from
different document subtrees. Future work could route between the two algorithms
based on query complexity classification.

\subsection{Offline Evaluation Limitations}
All experiments reported here use \textbf{offline (extractive) evaluation mode}:
traversal uses TF-IDF keyword matching in place of the full LLM scoring
$P(v\mid q)$, and answers are assembled by concatenating retrieved node summaries
rather than generating fluent text. This mode was necessary because the LLM
free-tier rate limit ($\approx$1 request per 13~s) renders large-scale online
evaluation prohibitively slow across $204+20+42=266$ questions $\times$ 6 systems.
The offline mode introduces two conservative biases: (1) \emph{traversal
quality}---keyword scoring is a weaker proxy for semantic relevance than LLM
scoring, penalizing both TreeRAG systems relative to their online performance;
and (2) \emph{answer quality}---extractive concatenation produces longer, less
precise answers than LLM-generated text, inflating context-length measurements
and deflating short-reference ROUGE-L scores for generation-heavy systems
(FlatRAG, RAPTOR). The relative \emph{ranking} of systems is expected to be
preserved in online evaluation; the \emph{absolute values} for BERTScore and
LLM-Judge will increase for TreeRAG systems. Online evaluation with real LLM
generation is planned for the camera-ready version.

\subsection{Limitations}
\textbf{LLM API dependency.} Indexing requires an LLM API call per document; rate
limits on free tiers slow bulk indexing. \textbf{Structural prompt brittleness.}
Documents with inconsistent heading styles (e.g.\ scanned PDFs with OCR
artifacts) may yield malformed index trees requiring manual correction.
\textbf{DSPy optimization.} We attempted learned node scoring via
DSPy~\cite{khattab2024dspy} with an open 70B model; optimization yielded no
improvement (0.0\% gain), likely due to insufficient training signal ($<100$
labeled examples), and we exclude this component from the main system.
\textbf{Evaluation dataset size.} The medical benchmark ($n=42$) and HotpotQA
subset ($n=20$) are small; larger-scale evaluation is needed to confirm findings.

\section{Conclusion}
\label{sec:conclusion}
We presented TreeRAG, a hierarchical document RAG system that indexes PDFs into
page-referenced JSON trees using zero-shot LLM prompting, traverses them with DFS
or Beam Search, compresses retrieved context, and validates generated answers
with a lightweight five-signal hallucination detector. On a 204-question general
benchmark, TreeRAG-DFS achieves ROUGE-L 0.451 ($+39.5\%$ over BM25, $+81.9\%$
over FlatRAG, $+132.4\%$ over RAPTOR; all $p<0.001$, $d=0.483$--$1.205$). On a
20-question multi-hop HotpotQA subset, TreeRAG-Beam achieves ROUGE-L 0.169
($+84.4\%$ over BM25, $p<0.001$, $d=1.358$). TreeRAG-DFS uses 37.5\% fewer
context tokens than BM25 while Pareto-dominating every baseline on the
accuracy--context trade-off. On the 42-question medical benchmark, TreeRAG-DFS
achieves ROUGE-L 0.366 with 100\% medical entity recall and full page-citation
availability---outperforming RAPTOR (ROUGE-L 0.053, entity recall 89.5\%, 0\%
citations). The system provides page-level citations, domain-adaptive prompting
for five domains, multilingual support, and a production-ready API with caching,
rate limiting, and async task queuing. Future work will pursue (1) online
LLM-evaluated benchmarks at scale; (2) Graph RAG integration~\cite{edge2025graphrag}
for inter-document relationship modeling; (3) multi-modal indexing for figures
and tables; and (4) active learning for traversal policy optimization from user
feedback signals.

%% --------------------------------------------------------------------
%% GenAI Usage Disclosure (required by CIKM 2026; remove for venues that
%% do not request it). Edit to reflect actual usage.
%% --------------------------------------------------------------------
\section*{GenAI Usage Disclosure}
Generative AI was used as a core component of the described system (LLM-driven
tree indexing, answer generation, and LLM-as-judge evaluation), as detailed in
Sections~\ref{sec:method}--\ref{sec:experiments}. In preparing this manuscript,
generative AI tools were also used for language editing and reference
formatting; all technical content, experimental design, analysis, and claims
were produced and verified by the authors.

\begin{thebibliography}{99}

\bibitem{asai2024selfrag}
A.~Asai, Z.~Wu, Y.~Wang, A.~Sil, and H.~Hajishirzi. 2024.
Self-RAG: Learning to Retrieve, Generate, and Critique through Self-Reflection.
In \emph{Proc.\ ICLR 2024}. arXiv:2310.11511.

\bibitem{brown2020gpt3}
T.~B. Brown et al. 2020. Language Models are Few-Shot Learners.
In \emph{NeurIPS 2020}, vol.~33, 1877--1901. arXiv:2005.14165.

\bibitem{edge2025graphrag}
D.~Edge, H.~Trinh, N.~Cheng, J.~Bradley, A.~Chao, A.~Mody, S.~Truitt,
D.~Metropolitansky, R.~O. Ness, and J.~Larson. 2025.
From Local to Global: A GraphRAG Approach to Query-Focused Summarization.
arXiv:2404.16130.

\bibitem{es2023ragas}
S.~Es, J.~James, L.~Espinosa-Anke, and S.~Schockaert. 2023.
RAGAS: Automated Evaluation of Retrieval Augmented Generation. arXiv:2309.15217.

\bibitem{formal2021splade}
T.~Formal, B.~Piwowarski, and S.~Clinchant. 2021. SPLADE: Sparse Lexical and
Expansion Model for First Stage Ranking. In \emph{Proc.\ SIGIR 2021}, 2288--2292.
doi:10.1145/3404835.3463098.

\bibitem{formal2024splate}
T.~Formal, S.~Clinchant, H.~D\'ejean, and C.~Lassance. 2024.
SPLATE: Sparse Late Interaction Retrieval. In \emph{Proc.\ SIGIR 2024}.
arXiv:2404.13950.

\bibitem{gao2024survey}
Y.~Gao, Y.~Xiong, X.~Gao, K.~Jia, J.~Pan, Y.~Bi, Y.~Dai, J.~Sun, M.~Wang, and
H.~Wang. 2024. Retrieval-Augmented Generation for Large Language Models: A
Survey. arXiv:2312.10997.

\bibitem{gupta2024survey}
S.~Gupta, R.~Ranjan, and S.~N. Singh. 2024. A Comprehensive Survey of
Retrieval-Augmented Generation (RAG): Evolution, Current Landscape and Future
Directions. arXiv:2410.12837.

\bibitem{huang2025survey}
L.~Huang et al. 2025. A Survey on Hallucination in Large Language Models:
Principles, Taxonomy, Challenges, and Open Questions.
\emph{ACM Trans.\ Inf.\ Syst.} 43(2), Article 42. doi:10.1145/3703155.

\bibitem{huang2025hirag}
H.~Huang, Y.~Huang, J.~Yang, Z.~Pan, Y.~Chen, K.~Ma, H.~Chen, and J.~Cheng. 2025.
Retrieval-Augmented Generation with Hierarchical Knowledge. arXiv:2503.10150.

\bibitem{jiang2023llmlingua}
H.~Jiang, Q.~Wu, C.-Y. Lin, Y.~Yang, and L.~Qiu. 2023. LLMLingua: Compressing
Prompts for Accelerated Inference of Large Language Models. In \emph{Proc.\
EMNLP 2023}, 13358--13376. arXiv:2310.05736.

\bibitem{johnson2019faiss}
J.~Johnson, M.~Douze, and H.~J\'egou. 2019. Billion-scale Similarity Search with
GPUs. \emph{IEEE Trans.\ Big Data} 7(3), 535--547. arXiv:1702.08734.

\bibitem{karpukhin2020dpr}
V.~Karpukhin, B.~O\u{g}uz, S.~Min, P.~Lewis, L.~Wu, S.~Edunov, D.~Chen, and
W.~Yih. 2020. Dense Passage Retrieval for Open-Domain Question Answering.
In \emph{Proc.\ EMNLP 2020}, 6769--6781.

\bibitem{khattab2020colbert}
O.~Khattab and M.~Zaharia. 2020. ColBERT: Efficient and Effective Passage Search
via Contextualized Late Interaction over BERT. In \emph{Proc.\ SIGIR 2020},
39--48. doi:10.1145/3397271.3401075.

\bibitem{khattab2024dspy}
O.~Khattab, A.~Singhvi, P.~Maheshwari, Z.~Zhang, K.~Santhanam, S.~Vardhamanan,
S.~Haq, A.~Sharma, T.~T. Joshi, H.~Moazam, H.~Miller, M.~Zaharia, and C.~Potts.
2024. DSPy: Compiling Declarative Language Model Calls into Self-Improving
Pipelines. In \emph{Proc.\ ICLR 2024}. arXiv:2310.03714.

\bibitem{kojima2022zero}
T.~Kojima, S.~S. Gu, M.~Reid, Y.~Matsuo, and Y.~Iwasawa. 2022. Large Language
Models are Zero-Shot Reasoners. In \emph{NeurIPS 2022}, vol.~35. arXiv:2205.11916.

\bibitem{kryscinski2020factcc}
W.~Kry\'sci\'nski, B.~McCann, C.~Xiong, and R.~Socher. 2020. Evaluating the
Factual Consistency of Abstractive Text Summarization. In \emph{Proc.\ EMNLP
2020}, 9332--9346. arXiv:1910.12840.

\bibitem{kwon2023paged}
W.~Kwon, Z.~Li, S.~Zhuang, Y.~Sheng, L.~Zheng, C.~H. Yu, J.~E. Gonzalez,
H.~Zhang, and I.~Stoica. 2023. Efficient Memory Management for Large Language
Model Serving with PagedAttention. In \emph{Proc.\ SOSP 2023}, 611--626.
doi:10.1145/3600006.3613165.

\bibitem{lewis2020rag}
P.~Lewis et al. 2020. Retrieval-Augmented Generation for Knowledge-Intensive NLP
Tasks. In \emph{NeurIPS 2020}, vol.~33, 9459--9474. arXiv:2005.11401.

\bibitem{lin2004rouge}
C.-Y. Lin. 2004. ROUGE: A Package for Automatic Evaluation of Summaries. In
\emph{Proc.\ ACL Workshop on Text Summarization Branches Out}, 74--81.

\bibitem{liu2022llamaindex}
J.~Liu. 2022. LlamaIndex. GitHub.

\bibitem{liu2024lostmiddle}
N.~F. Liu, K.~Lin, J.~Hewitt, A.~Paranjape, M.~Bevilacqua, F.~Petroni, and
P.~Liang. 2024. Lost in the Middle: How Language Models Use Long Contexts.
\emph{TACL} 12, 157--173. arXiv:2307.03172.

\bibitem{ma2021replication}
X.~Ma, R.~Sun, R.~Pradeep, and J.~Lin. 2021. A Replication Study of Dense Passage
Retrieval for Open-Domain Question Answering. arXiv:2104.05740.

\bibitem{malin2025review}
B.~Malin, T.~Kalganova, and N.~Boulgouris. 2025. A Review of Faithfulness Metrics
for Hallucination Assessment in Large Language Models. arXiv:2501.00269.

\bibitem{min2023factscore}
S.~Min, K.~Krishna, X.~Lyu, M.~Lewis, W.~Yih, P.~W. Koh, M.~Iyyer, L.~Zettlemoyer,
and H.~Hajishirzi. 2023. FActScore: Fine-grained Atomic Evaluation of Factual
Precision in Long Form Text Generation. In \emph{Proc.\ EMNLP 2023}.
arXiv:2305.14251.

\bibitem{nan2021entity}
F.~Nan, R.~Nallapati, Z.~Wang, C.~N. dos Santos, H.~Zhu, D.~Zhang, K.~McKeown,
and B.~Xiang. 2021. Entity-level Factual Consistency of Abstractive Text
Summarization. In \emph{Proc.\ EACL 2021}, 2727--2733. arXiv:2102.09130.

\bibitem{reimers2019sbert}
N.~Reimers and I.~Gurevych. 2019. Sentence-BERT: Sentence Embeddings using
Siamese BERT-Networks. In \emph{Proc.\ EMNLP 2019}, 3982--3992. arXiv:1908.10084.

\bibitem{robertson2009bm25}
S.~Robertson and H.~Zaragoza. 2009. The Probabilistic Relevance Framework: BM25
and Beyond. \emph{Found.\ Trends Inf.\ Retr.} 3(4), 333--389.

\bibitem{santhanam2022colbertv2}
K.~Santhanam, O.~Khattab, J.~Saad-Falcon, C.~Potts, and M.~Zaharia. 2022.
ColBERTv2: Effective and Efficient Retrieval via Lightweight Late Interaction.
In \emph{Proc.\ NAACL-HLT 2022}, 3715--3734. arXiv:2112.01488.

\bibitem{sarthi2024raptor}
P.~Sarthi, S.~Abdullah, A.~Tuli, S.~Khanna, A.~Goldie, and C.~D. Manning. 2024.
RAPTOR: Recursive Abstractive Processing for Tree-Organized Retrieval. In
\emph{Proc.\ ICLR 2024}. arXiv:2401.18059.

\bibitem{shuster2021retrieval}
K.~Shuster, S.~Poff, M.~Chen, D.~Kiela, and J.~Weston. 2021. Retrieval
Augmentation Reduces Hallucination in Conversation. In \emph{Findings of EMNLP
2021}, 3784--3803. arXiv:2104.07567.

\bibitem{trivedi2023ircot}
H.~Trivedi, N.~Balasubramanian, T.~Khot, and A.~Sabharwal. 2023. Interleaving
Retrieval with Chain-of-Thought Reasoning for Knowledge-Intensive Multi-Step
Questions. In \emph{Proc.\ ACL 2023}, 10014--10037. arXiv:2212.10509.

\bibitem{vectifyai2023pageindex}
Vectify AI. 2023. PageIndex: Reasoning-Based, Vectorless RAG via Document Tree
Search. GitHub.

\bibitem{wei2022cot}
J.~Wei, X.~Wang, D.~Schuurmans, M.~Bosma, B.~Ichter, F.~Xia, E.~Chi, Q.~Le, and
D.~Zhou. 2022. Chain-of-Thought Prompting Elicits Reasoning in Large Language
Models. In \emph{NeurIPS 2022}, vol.~35. arXiv:2201.11903.

\bibitem{yang2018hotpotqa}
Z.~Yang, P.~Qi, S.~Zhang, Y.~Bengio, W.~Cohen, R.~Salakhutdinov, and C.~D.
Manning. 2018. HotpotQA: A Dataset for Diverse, Explainable Multi-hop Question
Answering. In \emph{Proc.\ EMNLP 2018}, 2369--2380. arXiv:1809.09600.

\bibitem{zhang2020bertscore}
T.~Zhang, V.~Kishore, F.~Wu, K.~Q. Weinberger, and Y.~Artzi. 2020. BERTScore:
Evaluating Text Generation with BERT. In \emph{Proc.\ ICLR 2020}. arXiv:1904.09675.

\bibitem{zhang2024hierarchical}
X.~Zhang, M.~Wang, X.~Yang, D.~Wang, S.~Feng, and Y.~Zhang. 2024. Hierarchical
Retrieval-Augmented Generation Model with Rethink for Multi-hop Question
Answering. arXiv:2408.11875.

\bibitem{zheng2023judging}
L.~Zheng et al. 2023. Judging LLM-as-a-Judge with MT-Bench and Chatbot Arena.
In \emph{NeurIPS 2023}, Datasets and Benchmarks Track. arXiv:2306.05685.

\end{thebibliography}

%% --------------------------------------------------------------------
%% FIGURES. acmart places these floats; uncomment \includegraphics once
%% the image files are in the project. Filenames match the repo assets.
%% --------------------------------------------------------------------
\begin{figure}[t]
\centering
% \includegraphics[width=\columnwidth]{figure_1_architecture.png}
\caption{TreeRAG system pipeline: indexing (PDF $\rightarrow$ LLM $\rightarrow$
JSON tree) and reasoning (route $\rightarrow$ traverse $\rightarrow$ compress
$\rightarrow$ generate $\rightarrow$ verify).}
\Description{TreeRAG architecture diagram.}
\label{fig:arch}
\end{figure}

\begin{figure}[t]
\centering
% \includegraphics[width=\columnwidth]{figure_6_hallucination.png}
\caption{Five-signal hallucination detection pipeline scoring each answer
sentence against retrieved nodes.}
\Description{Hallucination detection illustration.}
\label{fig:halluc}
\end{figure}

\end{document}
